%% Sample Exam 3 -- Medical Image Analysis (ETH Zurich, 2026)
%%
%% Written in the style of Example_Exam_with_Answers.pdf.
%%
%% Compile the version WITH solutions:      pdflatex "\def\withanswers{}%% Sample Exam 3 -- Medical Image Analysis (ETH Zurich, 2026)
%%
%% Written in the style of Example_Exam_with_Answers.pdf.
%%
%% Compile the version WITH solutions:      pdflatex "\def\withanswers{}%% Sample Exam 3 -- Medical Image Analysis (ETH Zurich, 2026)
%%
%% Written in the style of Example_Exam_with_Answers.pdf.
%%
%% Compile the version WITH solutions:      pdflatex "\def\withanswers{}%% Sample Exam 3 -- Medical Image Analysis (ETH Zurich, 2026)
%%
%% Written in the style of Example_Exam_with_Answers.pdf.
%%
%% Compile the version WITH solutions:      pdflatex "\def\withanswers{}\input{Sample_Exam_3.tex}"
%% Compile the blank version (no answers):  pdflatex Sample_Exam_3.tex
%%
\ifx\withanswers\undefined
  \documentclass[11pt]{exam}
\else
  \documentclass[11pt,answers]{exam}
\fi

\usepackage[margin=2.4cm]{geometry}
\usepackage{amsmath,amssymb}
\usepackage{array}
\usepackage[T1]{fontenc}

% ---------------------------------------------------------------- page style
\pagestyle{foot}
\firstpageheader{}{}{}
\firstpagefooter{}{}{}
\runningheader{}{}{}
\runningfooter{}{}{Page \thepage}

% -------------------------------------------------------------- multiplechoice
% Reproduce the "(A) (B) (C) ..." labelling of the original exam.
\renewcommand{\choicelabel}{$\bigcirc$~(\Alph{choice})}

% ---------------------------------------------------------- true/false macros
% \TFT -> the statement is TRUE, \TFF -> the statement is FALSE.
% In the answer version the correct word is set in capitals and bold.
\newcommand{\TFT}{\ifprintanswers\textbf{TRUE}\hspace{1.1em}False\else True\hspace{1.1em}False\fi}
\newcommand{\TFF}{True\hspace{1.1em}\ifprintanswers\textbf{FALSE}\else False\fi}

% ------------------------------------------------------------- writing space
% Reserves room to write in the blank version; collapses in the answer version.
\newcommand{\answerspace}[1]{\ifprintanswers\else\vspace{#1}\fi}

\begin{document}

% ------------------------------------------------------------------ title page
\begin{center}
  {\LARGE\bfseries Sample Exam}\\[2mm]
  {\Large Medical Image Analysis}\\[6mm]
\end{center}

\vspace{4mm}
\noindent Student name:\hrulefill

\vspace{6mm}
\noindent Immatriculation number:\hrulefill

\vspace{10mm}

\begin{center}
  {\large\bfseries Exam Questions}
\end{center}

\noindent\textbf{Note:} All the questions have equal weight.

\vspace{4mm}

\begin{questions}

% =============================================================== Question 1 ===
\question
In 2D acquisitions, 2D image slices are acquired in one pre-defined direction one after the
other, and the final 3D volume is created by stacking the 2D slices. The voxel sizes are given
as in-plane (two values, the pixel size within the imaging plane) and through-plane (one value,
the slice thickness). Recall that the field-of-view is the real-world size of the entire image,
i.e.\ image size $\times$ voxel size. Based on this please compute the following.

\begin{parts}

  \part A \emph{coronal} acquisition with voxel size in-plane $0.9 \times 0.9$~mm$^2$ with 260
  pixels in both directions and voxel size through-plane $3.0$~mm with 64 pixels is acquired.
  What is the field-of-view in the head-to-toe, left-to-right and anterior-to-posterior
  directions, respectively?

  \begin{solution}
    Answer: $234 \times 234 \times 192$~mm$^3$.

    The coronal plane contains the head-to-toe and the left-to-right directions, so both in-plane
    directions give $260 \times 0.9 = 234$~mm. The through-plane direction of a coronal
    acquisition is anterior-to-posterior: $64 \times 3.0 = 192$~mm.
  \end{solution}
\answerspace{2.2cm}

  \part The protocol has to be changed so that the anterior-to-posterior coverage is at least
  210~mm, while the slice thickness stays at $3.0$~mm. How many slices have to be acquired?

  \begin{solution}
    Answer: 70 slices.

    $210 / 3.0 = 70$. The 64 slices of the original protocol only cover 192~mm.
  \end{solution}
\answerspace{1.8cm}

  \part You want to resample the volume of part (a) to isotropic voxels of
  $1.0 \times 1.0 \times 1.0$~mm$^3$. What are the resize factors and the resulting image size?
  In which of the three directions is the resampling least justified by the information actually
  contained in the data, and why?

  \begin{solution}
    Answer: resize factors $0.9 \times 0.9 \times 3.0$, resulting image size
    $234 \times 234 \times 192$~voxels.

    The resize factor per direction is the ratio of the old to the new voxel size: $0.9/1.0 = 0.9$
    head-to-toe, $0.9/1.0 = 0.9$ left-to-right and $3.0/1.0 = 3.0$ anterior-to-posterior.
    Applying these: $260 \cdot 0.9 = 234$, $260 \cdot 0.9 = 234$, $64 \cdot 3.0 = 192$. The
    field-of-view is unchanged, as it must be.

    The anterior-to-posterior direction. It is the through-plane direction, and it is being
    upsampled by a factor of 3: the interpolation invents two new samples between every pair of
    acquired slices. No new information is created --- the effective resolution in that direction
    is still that of a 3~mm slice, even though the voxel grid now says 1~mm.
  \end{solution}
\answerspace{4.5cm}

\end{parts}

% =============================================================== Question 2 ===
\question
Histogram equalization improves the contrast of an image by finding a monotonically increasing
intensity mapping $i \mapsto f(i)$ that flattens the histogram. Let $p(i)$ denote the probability
density function of the intensities, $P(i) = \int_0^i p(j)\,dj$ the cumulative distribution
function, and $i_{max}$ the maximum intensity. Which of the two mappings below is the one that
flattens the histogram? Please circle the correct choice.

\vspace{2mm}
\begin{center}
  $\displaystyle f(i) = i_{max}\, p(i)$
  \hspace{2.6cm}
  $\displaystyle f(i) = i_{max}\, P(i) = i_{max} \int_0^i p(j)\, dj$
\end{center}
\vspace{2mm}

\begin{solution}
  Answer: Right one --- the scaled \emph{cumulative} distribution function.

  Changing variables gives
  $p(f(i)) = p(i)\,\frac{di}{df(i)} = p(i) \cdot \frac{1}{p(i)} \cdot \frac{1}{i_{max}}
  = \frac{1}{i_{max}}$, i.e.\ a uniform (flat) density. Note also that only the CDF is guaranteed
  to be monotonically increasing, which the mapping has to be in order not to reorder intensities;
  the PDF is not.
\end{solution}
\answerspace{1.5cm}

% =============================================================== Question 3 ===
\question
When a Markov random field prior is used, segmentation can be written either as a posterior
maximization or as an energy minimization. $I$ is the image, $c$ the labelling, $G(x)$ the
neighbours of pixel $x$, and $Z$ a normalization constant.

\vspace{2mm}
\begin{center}
  $\displaystyle \max_c\ \log \Big[ \prod_{x \in \Omega} p(I(x) \mid c(x)) \Big]
   + \log \Big[ \tfrac{1}{Z} \exp\{-E(c)\} \Big]$
\end{center}
\vspace{1mm}
\begin{center}
  $\displaystyle \min_c\ \underbrace{\sum_{x \in \Omega} g\big(I(x) \mid \theta_{c(x)}\big)}_{\text{unary term}}
   \ +\ \underbrace{\sum_{x \in \Omega} \sum_{y \in G(x)} d\big(c(x), c(y)\big)}_{\text{pairwise term}}$
\end{center}
\vspace{2mm}

\begin{parts}
  \part Please indicate which term corresponds to which term in the two formulations.

  \begin{solution}
    Answer:
    \begin{align*}
      \log \textstyle\prod_{x} p(I(x) \mid c(x))
        &\ \Longleftrightarrow\ \text{the unary term}
        &&\text{(likelihood / data fidelity)}\\[1mm]
      \log \tfrac{1}{Z}\exp\{-E(c)\}
        &\ \Longleftrightarrow\ \text{the pairwise term}
        &&\text{(prior / neighbourhood consistency)}
    \end{align*}

    The correspondence requires the data model to be exponential, i.e.\
    $p(I(x) \mid c(x)) \propto \exp\{-g(I(x)\mid\theta_{c(x)})\}$, so that taking the logarithm of
    the product turns it into the unary sum. The pairwise term is exactly $E(c)$, and $\log Z$
    does not depend on $c$ and therefore drops out of the maximization.
  \end{solution}
\answerspace{2.8cm}

  \part Which theorem guarantees that a prior satisfying the Markov property can be written as a
  Gibbs distribution $\frac{1}{Z}\exp\{-E(c)\}$ with a locally defined energy $E$?

  \begin{solution}
    Answer: the Hammersley--Clifford theorem. It states that any probability distribution
    satisfying a Markovian property is a Gibbs distribution for an appropriate locally defined
    energy, and vice versa.
  \end{solution}
\answerspace{1.8cm}
\end{parts}

% =============================================================== Question 4 ===
\question
Please indicate whether the statement is true or false by circling the correct choice in each of
the following.

\begin{parts}
  \part \TFT\ \ A rigid-body transformation in 3D has 6 degrees of freedom: three rotation angles
  and three translations.

  \part \TFT\ \ For \emph{isotropic} scaling, $SRx = RSx$, i.e.\ the order in which rotation and
  scaling are applied does not matter.

  \part \TFF\ \ In homogeneous coordinates, a 3D affine transformation is written as a
  $3 \times 3$ matrix.

  \part \TFF\ \ The determinant of the Jacobian of a transformation is larger than 1 in the
  regions that the transformation compresses.
\end{parts}

\begin{solution}
  (a) TRUE --- three rotations, one about each axis, plus a translation vector in $\mathbb{R}^3$.
  In 2D the rigid model has 3 degrees of freedom (one angle, two translations).

  (b) TRUE --- with $S = sI$ the scaling commutes with any matrix. This is no longer the case for
  \emph{anisotropic} scaling, where $S = \mathrm{diag}(s_1, s_2, s_3)$: there the order matters a
  great deal, and applying the factors in the wrong order introduces shearing.

  (c) FALSE --- it is a $4 \times 4$ matrix,
  $\tilde A = \left[\begin{smallmatrix} A & t \\ \mathbf{0} & 1 \end{smallmatrix}\right]$ with
  $A \in \mathbb{R}^{3\times3}$ and $t \in \mathbb{R}^3$. The $3 \times 3$ homogeneous matrix is
  the 2D case.

  (d) FALSE --- it is the other way around: $\det(J) > 1$ means local \emph{expansion},
  $\det(J) < 1$ local compression, and $\det(J) = 1$ no volume change. The divergence
  $\nabla \cdot T$, the trace of the Jacobian, carries related information.
\end{solution}

% =============================================================== Question 5 ===
\question
A free-form deformation (FFD) model with cubic B-splines is used for the non-linear registration
of two 3D images. The control points are placed on a uniform grid of $10 \times 12 \times 8$
points. How many parameters does this transformation model have, and how does this compare with
the number of parameters of a 3D affine transformation?

\begin{solution}
  Answer: $10 \cdot 12 \cdot 8 = 960$ control points, each carrying a displacement vector in
  $\mathbb{R}^3$, hence $960 \times 3 = 2880$ parameters. A 3D affine transformation has 12
  parameters (3 rotations $+$ 3 translations $+$ 3 scalings $+$ 3 shears).

  The FFD model therefore has roughly 240 times more degrees of freedom, which is exactly the
  point: the affine model is global and cannot capture local deformation, while the FFD model can.
  Note that the parameter count of the FFD depends only on the number of control points, not on
  the number of voxels --- that is what makes it tractable compared with one displacement vector
  per voxel.
\end{solution}
\answerspace{3.2cm}

% =============================================================== Question 6 ===
\question
The N3 algorithm corrects the bias field in MRI. It works in the log domain, where the observed
image is $j(x) = b(x) + i(x)$ with $b$ the bias field and $i$ the bias-free image, and where the
bias field is assumed to be independent of the image intensities. Each iteration alternates
between two estimation problems. Please briefly describe what each of them does, and explain why
the algorithm has to iterate with narrow Gaussians instead of correcting the bias in a single
step. You do not need to write down the equations.

\begin{solution}
  Answer: \emph{Estimating the distribution of the bias-free intensities.} Because $b$ and $i$
  are independent and add in the log domain, the observed intensity distribution is the
  convolution $J = I * B$. Since the multiplicative bias is close to 1, its logarithm is close to
  0, so $B$ can be approximated by a Gaussian with a small standard deviation and $J$ is a blurred
  version of $I$. The distribution $I$ is therefore recovered by \emph{sharpening} $J$, using a
  Wiener-style deblurring filter in the Fourier domain.

  \emph{Estimating the field.} Given $I$ and $B$, the expected bias-free intensity corresponding
  to an observed intensity $j$ is $E[i \mid j]$, so the bias estimate is $b_e(j) = j - E[i \mid j]$.
  This is a one-dimensional integral and is evaluated numerically. It is computed at anchor points
  and then smoothed over space with splines or radial basis functions --- this is where the prior
  knowledge that the bias field is spatially smooth enters --- and subtracted from $j$.

  Why iterate: $B$ is not actually known, it is only assumed to be Gaussian. Deconvolving with a
  wide Gaussian is unstable and can produce estimates that are not valid distributions at all.
  Since a wide Gaussian is a convolution of narrow ones, the correction is instead applied as a
  sequence of small, stable steps. Iterating also compensates for the other approximations made
  along the way: the noise term was ignored, and the deblurring filter is itself only approximate.
\end{solution}
\answerspace{7.5cm}

% =============================================================== Question 7 ===
\question
Which of the following statements on the loss metrics used for intensity-based registration is
\textbf{inaccurate}?

\begin{choices}
  \choice The sum of squared differences assumes that the intensities of the two images are
  directly comparable, which restricts it to intra-modality registration.
  \choice Normalized cross-correlation attains its maximum magnitude when the two images are
  related by a linear intensity relationship that is the same at every pixel; both positive and
  negative correlations indicate alignment, which is why the squared value is used in the loss.
  \choice Mutual information is maximized when the joint intensity distribution equals the product
  of the two marginals, that is, when the intensities of the two images are statistically
  independent.
  \choice Mutual information can capture more complicated, non-linear relationships between the
  intensities than normalized cross-correlation, but it is harder to compute.
\end{choices}

\begin{solution}
  Answer: C.

  It states exactly the condition under which mutual information is \emph{zero}:
  $\mathrm{MI} = 0$ if and only if $p(I,J) = p(I)p(J)$, which means the two images tell us nothing
  about each other. Registration \emph{maximizes} MI, because aligned images have a strong
  statistical relationship and their joint distribution is then as far as possible from the
  product of the marginals.
\end{solution}

% =============================================================== Question 8 ===
\question
Which of the following statements about interpretability and explainable AI (XAI) in medical image
analysis is \textbf{incorrect}?

\begin{choices}
  \choice Black-box, or model-agnostic, methods such as LIME require no access to the internals of
  the model, whereas gradient-based methods such as Grad-CAM are white-box methods.
  \choice Grad-CAM generalizes CAM: the importance weight of a feature map is the global average of
  the gradients of the class score with respect to that feature map, which makes it applicable to
  any CNN, while CAM requires a particular architecture.
  \choice Because saliency maps show which parts of the image the model actually uses, a model whose
  saliency lies largely outside the anatomical region of interest is necessarily performing poorly
  on the task.
  \choice TCAV explains a prediction in terms of human-defined concepts by measuring the directional
  derivative of the class score along a concept activation vector, and can be used to spot biases in
  a data set.
\end{choices}

\begin{solution}
  Answer: C.

  High accuracy and misplaced attention are entirely compatible --- that is precisely what makes
  shortcut learning dangerous. In the glioma grading case study, normalizing the intensities over
  the whole image instead of over the brain area \emph{raised} the AUC from 0.8857 to 0.9841, while
  the saliency maps showed the model attending outside the brain: the model was exploiting an
  artefact of the preprocessing. The pneumonia classifier that learned to recognize the clinical
  site from hospital-specific markers is the same story. Saliency outside the region of interest is
  a reason to distrust a good score, not evidence of a bad one.
\end{solution}

% =============================================================== Question 9 ===
\question
Which of the following statements about self-attention, as it was built up in the lecture, is
\textbf{incorrect}?

\begin{choices}
  \choice Parameter-free dot-product attention, $Y = \mathrm{Softmax}(XX^T)X$, has no capacity to
  learn from data, and every feature within a token plays an equal role in determining the attention
  coefficients.
  \choice Introducing a single learnable matrix $U$ and computing
  $Y = \mathrm{Softmax}(XUU^TX^T)\,XU$ still yields a symmetric matrix of attention scores, which is
  why separate key and query projections are introduced.
  \choice The scaling by $1/\sqrt{D_K}$ is used because the gradients of the softmax become very
  small for very large or very small inputs; if the entries of $Q$ and $K$ have unit variance, the
  entries of $QK^T$ have variance $D_K$.
  \choice In multi-head self-attention the outputs of the individual heads are averaged, so that the
  output dimension equals $D_V$ irrespective of the number of heads.
\end{choices}

\begin{solution}
  Answer: D.

  The heads are \emph{concatenated}, not averaged, and the concatenation is then projected back:
  $Y = \mathrm{Concat}(H_1, \dots, H_H)\,W_0$ with
  $\mathrm{Concat}(\cdot) \in \mathbb{R}^{N \times H D_V}$ and $W_0 \in \mathbb{R}^{H D_V \times D}$.
  Averaging would defeat the purpose of using several heads, which --- by analogy with the multiple
  filters of a convolutional layer --- is to let different heads attend to different patterns.
\end{solution}

% ============================================================== Question 10 ===
\question
In the Noise2Noise setting you have no clean images at all. For each input $x_n$ you only have
several \emph{noisy} versions $y_n^m$ of the underlying clean image $y_n$, and you train with
$\arg\min_\theta \sum_n \sum_m \|y_n^m - f(x_n;\theta)\|^2$. Under which condition does this give
the same solution as supervised training on clean targets, in the limit of infinite data?

\begin{choices}
  \choice When the noise is multiplicative rather than additive.
  \choice When the average of the noisy targets is the clean image, $\mathbb{E}[\hat y \mid x] = y$,
  because the prediction that minimizes an $\ell_2$ loss is the mean of the targets.
  \choice When at least half of the training targets are clean images.
  \choice When an adversarial loss is added to the $\ell_2$ loss, so that the outputs are forced to
  look realistic.
  \choice Never --- training on noisy targets always biases the solution towards the noise.
\end{choices}

\begin{solution}
  Answer: B.

  Minimizing $\mathbb{E}_{\hat y}[\|f(x;\theta) - \hat y\|^2]$ over the prediction gives
  $f(x;\theta) = \mathbb{E}_{\hat y}[\hat y]$: the $\ell_2$-optimal prediction is the mean of the
  targets it is trained against. So if the noisy targets are distributed around the clean image,
  $\mathbb{E}[\hat y \mid x] = y$, the network converges to the clean image even though it never
  sees one. The network cannot fit the individual noise realizations, because a single set of
  parameters has to satisfy all of them at once, and only their mean does.
\end{solution}

% ============================================================== Question 11 ===
\question
Consider a classification network that takes grey-scale input images of size $64 \times 64$ with
one input channel. All convolutions are \emph{valid} (no padding, $P = 0$), and the network uses
strided convolutions instead of pooling:

\begin{itemize}
  \item \textbf{Conv1}: 16 kernels with kernel size $K = 5 \times 5$ and stride $S = 2$ in each
    direction.
  \item \textbf{Conv2}: 32 kernels with $K = 5 \times 5$ and $S = 2$.
  \item \textbf{Conv3}: 64 kernels with $K = 3 \times 3$ and $S = 2$.
\end{itemize}

\noindent Recall the channel size of a strided, valid convolution:
$N_l = \lceil (N_{l-1} - k_1)/s_1 \rceil$. Circle True or False:

\begin{parts}
  \part \TFT\ \ The output of Conv3 has a spatial size of $5 \times 5$.

  \part \TFF\ \ The global receptive field of a neuron in Conv3 is $13 \times 13$.

  \part \TFF\ \ The global stride of a neuron in Conv3 is $6 \times 6$, since the three strides of
  2 add up.

  \part \TFT\ \ The total number of convolutional kernel weights (excluding the biases) is
  $5 \cdot 5 \cdot 1 \cdot 16 + 5 \cdot 5 \cdot 16 \cdot 32 + 3 \cdot 3 \cdot 32 \cdot 64
  = 31\,632$.
\end{parts}

\begin{solution}
  (a) TRUE --- applying $N_l = \lceil (N_{l-1} - k)/s \rceil$ layer by layer:
  Conv1 gives $\lceil (64-5)/2 \rceil = \lceil 29.5 \rceil = 30$;
  Conv2 gives $\lceil (30-5)/2 \rceil = \lceil 12.5 \rceil = 13$;
  Conv3 gives $\lceil (13-3)/2 \rceil = 5$.

  (b) FALSE --- it is $21 \times 21$. Tracking the receptive field $R$ and the jump $j$ from
  $R = 1$, $j = 1$ with $R \leftarrow R + (K-1)\cdot j$ and $j \leftarrow j \cdot S$:
  Conv1 $\rightarrow R = 5,\ j = 2$; Conv2 $\rightarrow R = 13,\ j = 4$;
  Conv3 $\rightarrow R = 21,\ j = 8$. The value $13$ is the receptive field one layer earlier,
  after Conv2.

  (c) FALSE --- strides compound multiplicatively: $2 \cdot 2 \cdot 2 = 8$, so the global stride
  is $8 \times 8$.

  (d) TRUE --- $400 + 12\,800 + 18\,432 = 31\,632$. Note that strided convolutions reduce the
  spatial dimensions quickly, but unlike max pooling they do \emph{not} buy any translation
  invariance, and they discard information --- the more so the larger the stride.
\end{solution}

% ============================================================== Question 12 ===
\question
Which of the following statements on restoration, synthesis and self-supervised learning are
correct (multiple answers are possible)?

\begin{choices}
  \choice Losses of the MSE type tend to produce over-smoothed restorations, which is one
  motivation for perceptual losses, where the two images are passed through a separately trained
  network and their intermediate features are compared.
  \choice In CycleGAN-style synthesis from unpaired data, the cycle-consistency term requires that
  mapping a source image to the target domain and back returns the original image.
  \choice In contrastive learning, the two views of an image have to be produced by geometric
  transformations only, since intensity changes would destroy the surrogate label.
  \choice A masked autoencoder uses a larger encoder on the visible parts of the image and a
  smaller decoder to reconstruct the masked parts; the trained encoder is then used as a backbone
  for downstream tasks.
  \choice Adversarial losses compare individual pairs of images, which is what makes them suitable
  for unpaired data sets.
\end{choices}

\begin{solution}
  Answer: (A), (B) and (D) are correct.

  (C) is wrong: the transformations used in contrastive learning are explicitly \emph{not}
  restricted to geometric ones and routinely include intensity changes --- the whole point is that
  the representation should be invariant to them. (The one place where transformations do have to
  be handled carefully is the \emph{local} contrastive loss, where they must not destroy the patch
  correspondences, so they are either restricted or undone before the patches are compared.)

  (E) is wrong in its premise, and the premise is what makes the conclusion work. An adversarial
  loss is a \emph{distributional} distance: the discriminator is not shown a pair of corresponding
  images, it is asked whether a single image looks like it came from the real set. That is exactly
  why it can be used when no correspondence between source and target images exists.
\end{solution}

\end{questions}

\end{document}
"
%% Compile the blank version (no answers):  pdflatex Sample_Exam_3.tex
%%
\ifx\withanswers\undefined
  \documentclass[11pt]{exam}
\else
  \documentclass[11pt,answers]{exam}
\fi

\usepackage[margin=2.4cm]{geometry}
\usepackage{amsmath,amssymb}
\usepackage{array}
\usepackage[T1]{fontenc}

% ---------------------------------------------------------------- page style
\pagestyle{foot}
\firstpageheader{}{}{}
\firstpagefooter{}{}{}
\runningheader{}{}{}
\runningfooter{}{}{Page \thepage}

% -------------------------------------------------------------- multiplechoice
% Reproduce the "(A) (B) (C) ..." labelling of the original exam.
\renewcommand{\choicelabel}{$\bigcirc$~(\Alph{choice})}

% ---------------------------------------------------------- true/false macros
% \TFT -> the statement is TRUE, \TFF -> the statement is FALSE.
% In the answer version the correct word is set in capitals and bold.
\newcommand{\TFT}{\ifprintanswers\textbf{TRUE}\hspace{1.1em}False\else True\hspace{1.1em}False\fi}
\newcommand{\TFF}{True\hspace{1.1em}\ifprintanswers\textbf{FALSE}\else False\fi}

% ------------------------------------------------------------- writing space
% Reserves room to write in the blank version; collapses in the answer version.
\newcommand{\answerspace}[1]{\ifprintanswers\else\vspace{#1}\fi}

\begin{document}

% ------------------------------------------------------------------ title page
\begin{center}
  {\LARGE\bfseries Sample Exam}\\[2mm]
  {\Large Medical Image Analysis}\\[6mm]
\end{center}

\vspace{4mm}
\noindent Student name:\hrulefill

\vspace{6mm}
\noindent Immatriculation number:\hrulefill

\vspace{10mm}

\begin{center}
  {\large\bfseries Exam Questions}
\end{center}

\noindent\textbf{Note:} All the questions have equal weight.

\vspace{4mm}

\begin{questions}

% =============================================================== Question 1 ===
\question
In 2D acquisitions, 2D image slices are acquired in one pre-defined direction one after the
other, and the final 3D volume is created by stacking the 2D slices. The voxel sizes are given
as in-plane (two values, the pixel size within the imaging plane) and through-plane (one value,
the slice thickness). Recall that the field-of-view is the real-world size of the entire image,
i.e.\ image size $\times$ voxel size. Based on this please compute the following.

\begin{parts}

  \part A \emph{coronal} acquisition with voxel size in-plane $0.9 \times 0.9$~mm$^2$ with 260
  pixels in both directions and voxel size through-plane $3.0$~mm with 64 pixels is acquired.
  What is the field-of-view in the head-to-toe, left-to-right and anterior-to-posterior
  directions, respectively?

  \begin{solution}
    Answer: $234 \times 234 \times 192$~mm$^3$.

    The coronal plane contains the head-to-toe and the left-to-right directions, so both in-plane
    directions give $260 \times 0.9 = 234$~mm. The through-plane direction of a coronal
    acquisition is anterior-to-posterior: $64 \times 3.0 = 192$~mm.
  \end{solution}
\answerspace{2.2cm}

  \part The protocol has to be changed so that the anterior-to-posterior coverage is at least
  210~mm, while the slice thickness stays at $3.0$~mm. How many slices have to be acquired?

  \begin{solution}
    Answer: 70 slices.

    $210 / 3.0 = 70$. The 64 slices of the original protocol only cover 192~mm.
  \end{solution}
\answerspace{1.8cm}

  \part You want to resample the volume of part (a) to isotropic voxels of
  $1.0 \times 1.0 \times 1.0$~mm$^3$. What are the resize factors and the resulting image size?
  In which of the three directions is the resampling least justified by the information actually
  contained in the data, and why?

  \begin{solution}
    Answer: resize factors $0.9 \times 0.9 \times 3.0$, resulting image size
    $234 \times 234 \times 192$~voxels.

    The resize factor per direction is the ratio of the old to the new voxel size: $0.9/1.0 = 0.9$
    head-to-toe, $0.9/1.0 = 0.9$ left-to-right and $3.0/1.0 = 3.0$ anterior-to-posterior.
    Applying these: $260 \cdot 0.9 = 234$, $260 \cdot 0.9 = 234$, $64 \cdot 3.0 = 192$. The
    field-of-view is unchanged, as it must be.

    The anterior-to-posterior direction. It is the through-plane direction, and it is being
    upsampled by a factor of 3: the interpolation invents two new samples between every pair of
    acquired slices. No new information is created --- the effective resolution in that direction
    is still that of a 3~mm slice, even though the voxel grid now says 1~mm.
  \end{solution}
\answerspace{4.5cm}

\end{parts}

% =============================================================== Question 2 ===
\question
Histogram equalization improves the contrast of an image by finding a monotonically increasing
intensity mapping $i \mapsto f(i)$ that flattens the histogram. Let $p(i)$ denote the probability
density function of the intensities, $P(i) = \int_0^i p(j)\,dj$ the cumulative distribution
function, and $i_{max}$ the maximum intensity. Which of the two mappings below is the one that
flattens the histogram? Please circle the correct choice.

\vspace{2mm}
\begin{center}
  $\displaystyle f(i) = i_{max}\, p(i)$
  \hspace{2.6cm}
  $\displaystyle f(i) = i_{max}\, P(i) = i_{max} \int_0^i p(j)\, dj$
\end{center}
\vspace{2mm}

\begin{solution}
  Answer: Right one --- the scaled \emph{cumulative} distribution function.

  Changing variables gives
  $p(f(i)) = p(i)\,\frac{di}{df(i)} = p(i) \cdot \frac{1}{p(i)} \cdot \frac{1}{i_{max}}
  = \frac{1}{i_{max}}$, i.e.\ a uniform (flat) density. Note also that only the CDF is guaranteed
  to be monotonically increasing, which the mapping has to be in order not to reorder intensities;
  the PDF is not.
\end{solution}
\answerspace{1.5cm}

% =============================================================== Question 3 ===
\question
When a Markov random field prior is used, segmentation can be written either as a posterior
maximization or as an energy minimization. $I$ is the image, $c$ the labelling, $G(x)$ the
neighbours of pixel $x$, and $Z$ a normalization constant.

\vspace{2mm}
\begin{center}
  $\displaystyle \max_c\ \log \Big[ \prod_{x \in \Omega} p(I(x) \mid c(x)) \Big]
   + \log \Big[ \tfrac{1}{Z} \exp\{-E(c)\} \Big]$
\end{center}
\vspace{1mm}
\begin{center}
  $\displaystyle \min_c\ \underbrace{\sum_{x \in \Omega} g\big(I(x) \mid \theta_{c(x)}\big)}_{\text{unary term}}
   \ +\ \underbrace{\sum_{x \in \Omega} \sum_{y \in G(x)} d\big(c(x), c(y)\big)}_{\text{pairwise term}}$
\end{center}
\vspace{2mm}

\begin{parts}
  \part Please indicate which term corresponds to which term in the two formulations.

  \begin{solution}
    Answer:
    \begin{align*}
      \log \textstyle\prod_{x} p(I(x) \mid c(x))
        &\ \Longleftrightarrow\ \text{the unary term}
        &&\text{(likelihood / data fidelity)}\\[1mm]
      \log \tfrac{1}{Z}\exp\{-E(c)\}
        &\ \Longleftrightarrow\ \text{the pairwise term}
        &&\text{(prior / neighbourhood consistency)}
    \end{align*}

    The correspondence requires the data model to be exponential, i.e.\
    $p(I(x) \mid c(x)) \propto \exp\{-g(I(x)\mid\theta_{c(x)})\}$, so that taking the logarithm of
    the product turns it into the unary sum. The pairwise term is exactly $E(c)$, and $\log Z$
    does not depend on $c$ and therefore drops out of the maximization.
  \end{solution}
\answerspace{2.8cm}

  \part Which theorem guarantees that a prior satisfying the Markov property can be written as a
  Gibbs distribution $\frac{1}{Z}\exp\{-E(c)\}$ with a locally defined energy $E$?

  \begin{solution}
    Answer: the Hammersley--Clifford theorem. It states that any probability distribution
    satisfying a Markovian property is a Gibbs distribution for an appropriate locally defined
    energy, and vice versa.
  \end{solution}
\answerspace{1.8cm}
\end{parts}

% =============================================================== Question 4 ===
\question
Please indicate whether the statement is true or false by circling the correct choice in each of
the following.

\begin{parts}
  \part \TFT\ \ A rigid-body transformation in 3D has 6 degrees of freedom: three rotation angles
  and three translations.

  \part \TFT\ \ For \emph{isotropic} scaling, $SRx = RSx$, i.e.\ the order in which rotation and
  scaling are applied does not matter.

  \part \TFF\ \ In homogeneous coordinates, a 3D affine transformation is written as a
  $3 \times 3$ matrix.

  \part \TFF\ \ The determinant of the Jacobian of a transformation is larger than 1 in the
  regions that the transformation compresses.
\end{parts}

\begin{solution}
  (a) TRUE --- three rotations, one about each axis, plus a translation vector in $\mathbb{R}^3$.
  In 2D the rigid model has 3 degrees of freedom (one angle, two translations).

  (b) TRUE --- with $S = sI$ the scaling commutes with any matrix. This is no longer the case for
  \emph{anisotropic} scaling, where $S = \mathrm{diag}(s_1, s_2, s_3)$: there the order matters a
  great deal, and applying the factors in the wrong order introduces shearing.

  (c) FALSE --- it is a $4 \times 4$ matrix,
  $\tilde A = \left[\begin{smallmatrix} A & t \\ \mathbf{0} & 1 \end{smallmatrix}\right]$ with
  $A \in \mathbb{R}^{3\times3}$ and $t \in \mathbb{R}^3$. The $3 \times 3$ homogeneous matrix is
  the 2D case.

  (d) FALSE --- it is the other way around: $\det(J) > 1$ means local \emph{expansion},
  $\det(J) < 1$ local compression, and $\det(J) = 1$ no volume change. The divergence
  $\nabla \cdot T$, the trace of the Jacobian, carries related information.
\end{solution}

% =============================================================== Question 5 ===
\question
A free-form deformation (FFD) model with cubic B-splines is used for the non-linear registration
of two 3D images. The control points are placed on a uniform grid of $10 \times 12 \times 8$
points. How many parameters does this transformation model have, and how does this compare with
the number of parameters of a 3D affine transformation?

\begin{solution}
  Answer: $10 \cdot 12 \cdot 8 = 960$ control points, each carrying a displacement vector in
  $\mathbb{R}^3$, hence $960 \times 3 = 2880$ parameters. A 3D affine transformation has 12
  parameters (3 rotations $+$ 3 translations $+$ 3 scalings $+$ 3 shears).

  The FFD model therefore has roughly 240 times more degrees of freedom, which is exactly the
  point: the affine model is global and cannot capture local deformation, while the FFD model can.
  Note that the parameter count of the FFD depends only on the number of control points, not on
  the number of voxels --- that is what makes it tractable compared with one displacement vector
  per voxel.
\end{solution}
\answerspace{3.2cm}

% =============================================================== Question 6 ===
\question
The N3 algorithm corrects the bias field in MRI. It works in the log domain, where the observed
image is $j(x) = b(x) + i(x)$ with $b$ the bias field and $i$ the bias-free image, and where the
bias field is assumed to be independent of the image intensities. Each iteration alternates
between two estimation problems. Please briefly describe what each of them does, and explain why
the algorithm has to iterate with narrow Gaussians instead of correcting the bias in a single
step. You do not need to write down the equations.

\begin{solution}
  Answer: \emph{Estimating the distribution of the bias-free intensities.} Because $b$ and $i$
  are independent and add in the log domain, the observed intensity distribution is the
  convolution $J = I * B$. Since the multiplicative bias is close to 1, its logarithm is close to
  0, so $B$ can be approximated by a Gaussian with a small standard deviation and $J$ is a blurred
  version of $I$. The distribution $I$ is therefore recovered by \emph{sharpening} $J$, using a
  Wiener-style deblurring filter in the Fourier domain.

  \emph{Estimating the field.} Given $I$ and $B$, the expected bias-free intensity corresponding
  to an observed intensity $j$ is $E[i \mid j]$, so the bias estimate is $b_e(j) = j - E[i \mid j]$.
  This is a one-dimensional integral and is evaluated numerically. It is computed at anchor points
  and then smoothed over space with splines or radial basis functions --- this is where the prior
  knowledge that the bias field is spatially smooth enters --- and subtracted from $j$.

  Why iterate: $B$ is not actually known, it is only assumed to be Gaussian. Deconvolving with a
  wide Gaussian is unstable and can produce estimates that are not valid distributions at all.
  Since a wide Gaussian is a convolution of narrow ones, the correction is instead applied as a
  sequence of small, stable steps. Iterating also compensates for the other approximations made
  along the way: the noise term was ignored, and the deblurring filter is itself only approximate.
\end{solution}
\answerspace{7.5cm}

% =============================================================== Question 7 ===
\question
Which of the following statements on the loss metrics used for intensity-based registration is
\textbf{inaccurate}?

\begin{choices}
  \choice The sum of squared differences assumes that the intensities of the two images are
  directly comparable, which restricts it to intra-modality registration.
  \choice Normalized cross-correlation attains its maximum magnitude when the two images are
  related by a linear intensity relationship that is the same at every pixel; both positive and
  negative correlations indicate alignment, which is why the squared value is used in the loss.
  \choice Mutual information is maximized when the joint intensity distribution equals the product
  of the two marginals, that is, when the intensities of the two images are statistically
  independent.
  \choice Mutual information can capture more complicated, non-linear relationships between the
  intensities than normalized cross-correlation, but it is harder to compute.
\end{choices}

\begin{solution}
  Answer: C.

  It states exactly the condition under which mutual information is \emph{zero}:
  $\mathrm{MI} = 0$ if and only if $p(I,J) = p(I)p(J)$, which means the two images tell us nothing
  about each other. Registration \emph{maximizes} MI, because aligned images have a strong
  statistical relationship and their joint distribution is then as far as possible from the
  product of the marginals.
\end{solution}

% =============================================================== Question 8 ===
\question
Which of the following statements about interpretability and explainable AI (XAI) in medical image
analysis is \textbf{incorrect}?

\begin{choices}
  \choice Black-box, or model-agnostic, methods such as LIME require no access to the internals of
  the model, whereas gradient-based methods such as Grad-CAM are white-box methods.
  \choice Grad-CAM generalizes CAM: the importance weight of a feature map is the global average of
  the gradients of the class score with respect to that feature map, which makes it applicable to
  any CNN, while CAM requires a particular architecture.
  \choice Because saliency maps show which parts of the image the model actually uses, a model whose
  saliency lies largely outside the anatomical region of interest is necessarily performing poorly
  on the task.
  \choice TCAV explains a prediction in terms of human-defined concepts by measuring the directional
  derivative of the class score along a concept activation vector, and can be used to spot biases in
  a data set.
\end{choices}

\begin{solution}
  Answer: C.

  High accuracy and misplaced attention are entirely compatible --- that is precisely what makes
  shortcut learning dangerous. In the glioma grading case study, normalizing the intensities over
  the whole image instead of over the brain area \emph{raised} the AUC from 0.8857 to 0.9841, while
  the saliency maps showed the model attending outside the brain: the model was exploiting an
  artefact of the preprocessing. The pneumonia classifier that learned to recognize the clinical
  site from hospital-specific markers is the same story. Saliency outside the region of interest is
  a reason to distrust a good score, not evidence of a bad one.
\end{solution}

% =============================================================== Question 9 ===
\question
Which of the following statements about self-attention, as it was built up in the lecture, is
\textbf{incorrect}?

\begin{choices}
  \choice Parameter-free dot-product attention, $Y = \mathrm{Softmax}(XX^T)X$, has no capacity to
  learn from data, and every feature within a token plays an equal role in determining the attention
  coefficients.
  \choice Introducing a single learnable matrix $U$ and computing
  $Y = \mathrm{Softmax}(XUU^TX^T)\,XU$ still yields a symmetric matrix of attention scores, which is
  why separate key and query projections are introduced.
  \choice The scaling by $1/\sqrt{D_K}$ is used because the gradients of the softmax become very
  small for very large or very small inputs; if the entries of $Q$ and $K$ have unit variance, the
  entries of $QK^T$ have variance $D_K$.
  \choice In multi-head self-attention the outputs of the individual heads are averaged, so that the
  output dimension equals $D_V$ irrespective of the number of heads.
\end{choices}

\begin{solution}
  Answer: D.

  The heads are \emph{concatenated}, not averaged, and the concatenation is then projected back:
  $Y = \mathrm{Concat}(H_1, \dots, H_H)\,W_0$ with
  $\mathrm{Concat}(\cdot) \in \mathbb{R}^{N \times H D_V}$ and $W_0 \in \mathbb{R}^{H D_V \times D}$.
  Averaging would defeat the purpose of using several heads, which --- by analogy with the multiple
  filters of a convolutional layer --- is to let different heads attend to different patterns.
\end{solution}

% ============================================================== Question 10 ===
\question
In the Noise2Noise setting you have no clean images at all. For each input $x_n$ you only have
several \emph{noisy} versions $y_n^m$ of the underlying clean image $y_n$, and you train with
$\arg\min_\theta \sum_n \sum_m \|y_n^m - f(x_n;\theta)\|^2$. Under which condition does this give
the same solution as supervised training on clean targets, in the limit of infinite data?

\begin{choices}
  \choice When the noise is multiplicative rather than additive.
  \choice When the average of the noisy targets is the clean image, $\mathbb{E}[\hat y \mid x] = y$,
  because the prediction that minimizes an $\ell_2$ loss is the mean of the targets.
  \choice When at least half of the training targets are clean images.
  \choice When an adversarial loss is added to the $\ell_2$ loss, so that the outputs are forced to
  look realistic.
  \choice Never --- training on noisy targets always biases the solution towards the noise.
\end{choices}

\begin{solution}
  Answer: B.

  Minimizing $\mathbb{E}_{\hat y}[\|f(x;\theta) - \hat y\|^2]$ over the prediction gives
  $f(x;\theta) = \mathbb{E}_{\hat y}[\hat y]$: the $\ell_2$-optimal prediction is the mean of the
  targets it is trained against. So if the noisy targets are distributed around the clean image,
  $\mathbb{E}[\hat y \mid x] = y$, the network converges to the clean image even though it never
  sees one. The network cannot fit the individual noise realizations, because a single set of
  parameters has to satisfy all of them at once, and only their mean does.
\end{solution}

% ============================================================== Question 11 ===
\question
Consider a classification network that takes grey-scale input images of size $64 \times 64$ with
one input channel. All convolutions are \emph{valid} (no padding, $P = 0$), and the network uses
strided convolutions instead of pooling:

\begin{itemize}
  \item \textbf{Conv1}: 16 kernels with kernel size $K = 5 \times 5$ and stride $S = 2$ in each
    direction.
  \item \textbf{Conv2}: 32 kernels with $K = 5 \times 5$ and $S = 2$.
  \item \textbf{Conv3}: 64 kernels with $K = 3 \times 3$ and $S = 2$.
\end{itemize}

\noindent Recall the channel size of a strided, valid convolution:
$N_l = \lceil (N_{l-1} - k_1)/s_1 \rceil$. Circle True or False:

\begin{parts}
  \part \TFT\ \ The output of Conv3 has a spatial size of $5 \times 5$.

  \part \TFF\ \ The global receptive field of a neuron in Conv3 is $13 \times 13$.

  \part \TFF\ \ The global stride of a neuron in Conv3 is $6 \times 6$, since the three strides of
  2 add up.

  \part \TFT\ \ The total number of convolutional kernel weights (excluding the biases) is
  $5 \cdot 5 \cdot 1 \cdot 16 + 5 \cdot 5 \cdot 16 \cdot 32 + 3 \cdot 3 \cdot 32 \cdot 64
  = 31\,632$.
\end{parts}

\begin{solution}
  (a) TRUE --- applying $N_l = \lceil (N_{l-1} - k)/s \rceil$ layer by layer:
  Conv1 gives $\lceil (64-5)/2 \rceil = \lceil 29.5 \rceil = 30$;
  Conv2 gives $\lceil (30-5)/2 \rceil = \lceil 12.5 \rceil = 13$;
  Conv3 gives $\lceil (13-3)/2 \rceil = 5$.

  (b) FALSE --- it is $21 \times 21$. Tracking the receptive field $R$ and the jump $j$ from
  $R = 1$, $j = 1$ with $R \leftarrow R + (K-1)\cdot j$ and $j \leftarrow j \cdot S$:
  Conv1 $\rightarrow R = 5,\ j = 2$; Conv2 $\rightarrow R = 13,\ j = 4$;
  Conv3 $\rightarrow R = 21,\ j = 8$. The value $13$ is the receptive field one layer earlier,
  after Conv2.

  (c) FALSE --- strides compound multiplicatively: $2 \cdot 2 \cdot 2 = 8$, so the global stride
  is $8 \times 8$.

  (d) TRUE --- $400 + 12\,800 + 18\,432 = 31\,632$. Note that strided convolutions reduce the
  spatial dimensions quickly, but unlike max pooling they do \emph{not} buy any translation
  invariance, and they discard information --- the more so the larger the stride.
\end{solution}

% ============================================================== Question 12 ===
\question
Which of the following statements on restoration, synthesis and self-supervised learning are
correct (multiple answers are possible)?

\begin{choices}
  \choice Losses of the MSE type tend to produce over-smoothed restorations, which is one
  motivation for perceptual losses, where the two images are passed through a separately trained
  network and their intermediate features are compared.
  \choice In CycleGAN-style synthesis from unpaired data, the cycle-consistency term requires that
  mapping a source image to the target domain and back returns the original image.
  \choice In contrastive learning, the two views of an image have to be produced by geometric
  transformations only, since intensity changes would destroy the surrogate label.
  \choice A masked autoencoder uses a larger encoder on the visible parts of the image and a
  smaller decoder to reconstruct the masked parts; the trained encoder is then used as a backbone
  for downstream tasks.
  \choice Adversarial losses compare individual pairs of images, which is what makes them suitable
  for unpaired data sets.
\end{choices}

\begin{solution}
  Answer: (A), (B) and (D) are correct.

  (C) is wrong: the transformations used in contrastive learning are explicitly \emph{not}
  restricted to geometric ones and routinely include intensity changes --- the whole point is that
  the representation should be invariant to them. (The one place where transformations do have to
  be handled carefully is the \emph{local} contrastive loss, where they must not destroy the patch
  correspondences, so they are either restricted or undone before the patches are compared.)

  (E) is wrong in its premise, and the premise is what makes the conclusion work. An adversarial
  loss is a \emph{distributional} distance: the discriminator is not shown a pair of corresponding
  images, it is asked whether a single image looks like it came from the real set. That is exactly
  why it can be used when no correspondence between source and target images exists.
\end{solution}

\end{questions}

\end{document}
"
%% Compile the blank version (no answers):  pdflatex Sample_Exam_3.tex
%%
\ifx\withanswers\undefined
  \documentclass[11pt]{exam}
\else
  \documentclass[11pt,answers]{exam}
\fi

\usepackage[margin=2.4cm]{geometry}
\usepackage{amsmath,amssymb}
\usepackage{array}
\usepackage[T1]{fontenc}

% ---------------------------------------------------------------- page style
\pagestyle{foot}
\firstpageheader{}{}{}
\firstpagefooter{}{}{}
\runningheader{}{}{}
\runningfooter{}{}{Page \thepage}

% -------------------------------------------------------------- multiplechoice
% Reproduce the "(A) (B) (C) ..." labelling of the original exam.
\renewcommand{\choicelabel}{$\bigcirc$~(\Alph{choice})}

% ---------------------------------------------------------- true/false macros
% \TFT -> the statement is TRUE, \TFF -> the statement is FALSE.
% In the answer version the correct word is set in capitals and bold.
\newcommand{\TFT}{\ifprintanswers\textbf{TRUE}\hspace{1.1em}False\else True\hspace{1.1em}False\fi}
\newcommand{\TFF}{True\hspace{1.1em}\ifprintanswers\textbf{FALSE}\else False\fi}

% ------------------------------------------------------------- writing space
% Reserves room to write in the blank version; collapses in the answer version.
\newcommand{\answerspace}[1]{\ifprintanswers\else\vspace{#1}\fi}

\begin{document}

% ------------------------------------------------------------------ title page
\begin{center}
  {\LARGE\bfseries Sample Exam}\\[2mm]
  {\Large Medical Image Analysis}\\[6mm]
\end{center}

\vspace{4mm}
\noindent Student name:\hrulefill

\vspace{6mm}
\noindent Immatriculation number:\hrulefill

\vspace{10mm}

\begin{center}
  {\large\bfseries Exam Questions}
\end{center}

\noindent\textbf{Note:} All the questions have equal weight.

\vspace{4mm}

\begin{questions}

% =============================================================== Question 1 ===
\question
In 2D acquisitions, 2D image slices are acquired in one pre-defined direction one after the
other, and the final 3D volume is created by stacking the 2D slices. The voxel sizes are given
as in-plane (two values, the pixel size within the imaging plane) and through-plane (one value,
the slice thickness). Recall that the field-of-view is the real-world size of the entire image,
i.e.\ image size $\times$ voxel size. Based on this please compute the following.

\begin{parts}

  \part A \emph{coronal} acquisition with voxel size in-plane $0.9 \times 0.9$~mm$^2$ with 260
  pixels in both directions and voxel size through-plane $3.0$~mm with 64 pixels is acquired.
  What is the field-of-view in the head-to-toe, left-to-right and anterior-to-posterior
  directions, respectively?

  \begin{solution}
    Answer: $234 \times 234 \times 192$~mm$^3$.

    The coronal plane contains the head-to-toe and the left-to-right directions, so both in-plane
    directions give $260 \times 0.9 = 234$~mm. The through-plane direction of a coronal
    acquisition is anterior-to-posterior: $64 \times 3.0 = 192$~mm.
  \end{solution}
\answerspace{2.2cm}

  \part The protocol has to be changed so that the anterior-to-posterior coverage is at least
  210~mm, while the slice thickness stays at $3.0$~mm. How many slices have to be acquired?

  \begin{solution}
    Answer: 70 slices.

    $210 / 3.0 = 70$. The 64 slices of the original protocol only cover 192~mm.
  \end{solution}
\answerspace{1.8cm}

  \part You want to resample the volume of part (a) to isotropic voxels of
  $1.0 \times 1.0 \times 1.0$~mm$^3$. What are the resize factors and the resulting image size?
  In which of the three directions is the resampling least justified by the information actually
  contained in the data, and why?

  \begin{solution}
    Answer: resize factors $0.9 \times 0.9 \times 3.0$, resulting image size
    $234 \times 234 \times 192$~voxels.

    The resize factor per direction is the ratio of the old to the new voxel size: $0.9/1.0 = 0.9$
    head-to-toe, $0.9/1.0 = 0.9$ left-to-right and $3.0/1.0 = 3.0$ anterior-to-posterior.
    Applying these: $260 \cdot 0.9 = 234$, $260 \cdot 0.9 = 234$, $64 \cdot 3.0 = 192$. The
    field-of-view is unchanged, as it must be.

    The anterior-to-posterior direction. It is the through-plane direction, and it is being
    upsampled by a factor of 3: the interpolation invents two new samples between every pair of
    acquired slices. No new information is created --- the effective resolution in that direction
    is still that of a 3~mm slice, even though the voxel grid now says 1~mm.
  \end{solution}
\answerspace{4.5cm}

\end{parts}

% =============================================================== Question 2 ===
\question
Histogram equalization improves the contrast of an image by finding a monotonically increasing
intensity mapping $i \mapsto f(i)$ that flattens the histogram. Let $p(i)$ denote the probability
density function of the intensities, $P(i) = \int_0^i p(j)\,dj$ the cumulative distribution
function, and $i_{max}$ the maximum intensity. Which of the two mappings below is the one that
flattens the histogram? Please circle the correct choice.

\vspace{2mm}
\begin{center}
  $\displaystyle f(i) = i_{max}\, p(i)$
  \hspace{2.6cm}
  $\displaystyle f(i) = i_{max}\, P(i) = i_{max} \int_0^i p(j)\, dj$
\end{center}
\vspace{2mm}

\begin{solution}
  Answer: Right one --- the scaled \emph{cumulative} distribution function.

  Changing variables gives
  $p(f(i)) = p(i)\,\frac{di}{df(i)} = p(i) \cdot \frac{1}{p(i)} \cdot \frac{1}{i_{max}}
  = \frac{1}{i_{max}}$, i.e.\ a uniform (flat) density. Note also that only the CDF is guaranteed
  to be monotonically increasing, which the mapping has to be in order not to reorder intensities;
  the PDF is not.
\end{solution}
\answerspace{1.5cm}

% =============================================================== Question 3 ===
\question
When a Markov random field prior is used, segmentation can be written either as a posterior
maximization or as an energy minimization. $I$ is the image, $c$ the labelling, $G(x)$ the
neighbours of pixel $x$, and $Z$ a normalization constant.

\vspace{2mm}
\begin{center}
  $\displaystyle \max_c\ \log \Big[ \prod_{x \in \Omega} p(I(x) \mid c(x)) \Big]
   + \log \Big[ \tfrac{1}{Z} \exp\{-E(c)\} \Big]$
\end{center}
\vspace{1mm}
\begin{center}
  $\displaystyle \min_c\ \underbrace{\sum_{x \in \Omega} g\big(I(x) \mid \theta_{c(x)}\big)}_{\text{unary term}}
   \ +\ \underbrace{\sum_{x \in \Omega} \sum_{y \in G(x)} d\big(c(x), c(y)\big)}_{\text{pairwise term}}$
\end{center}
\vspace{2mm}

\begin{parts}
  \part Please indicate which term corresponds to which term in the two formulations.

  \begin{solution}
    Answer:
    \begin{align*}
      \log \textstyle\prod_{x} p(I(x) \mid c(x))
        &\ \Longleftrightarrow\ \text{the unary term}
        &&\text{(likelihood / data fidelity)}\\[1mm]
      \log \tfrac{1}{Z}\exp\{-E(c)\}
        &\ \Longleftrightarrow\ \text{the pairwise term}
        &&\text{(prior / neighbourhood consistency)}
    \end{align*}

    The correspondence requires the data model to be exponential, i.e.\
    $p(I(x) \mid c(x)) \propto \exp\{-g(I(x)\mid\theta_{c(x)})\}$, so that taking the logarithm of
    the product turns it into the unary sum. The pairwise term is exactly $E(c)$, and $\log Z$
    does not depend on $c$ and therefore drops out of the maximization.
  \end{solution}
\answerspace{2.8cm}

  \part Which theorem guarantees that a prior satisfying the Markov property can be written as a
  Gibbs distribution $\frac{1}{Z}\exp\{-E(c)\}$ with a locally defined energy $E$?

  \begin{solution}
    Answer: the Hammersley--Clifford theorem. It states that any probability distribution
    satisfying a Markovian property is a Gibbs distribution for an appropriate locally defined
    energy, and vice versa.
  \end{solution}
\answerspace{1.8cm}
\end{parts}

% =============================================================== Question 4 ===
\question
Please indicate whether the statement is true or false by circling the correct choice in each of
the following.

\begin{parts}
  \part \TFT\ \ A rigid-body transformation in 3D has 6 degrees of freedom: three rotation angles
  and three translations.

  \part \TFT\ \ For \emph{isotropic} scaling, $SRx = RSx$, i.e.\ the order in which rotation and
  scaling are applied does not matter.

  \part \TFF\ \ In homogeneous coordinates, a 3D affine transformation is written as a
  $3 \times 3$ matrix.

  \part \TFF\ \ The determinant of the Jacobian of a transformation is larger than 1 in the
  regions that the transformation compresses.
\end{parts}

\begin{solution}
  (a) TRUE --- three rotations, one about each axis, plus a translation vector in $\mathbb{R}^3$.
  In 2D the rigid model has 3 degrees of freedom (one angle, two translations).

  (b) TRUE --- with $S = sI$ the scaling commutes with any matrix. This is no longer the case for
  \emph{anisotropic} scaling, where $S = \mathrm{diag}(s_1, s_2, s_3)$: there the order matters a
  great deal, and applying the factors in the wrong order introduces shearing.

  (c) FALSE --- it is a $4 \times 4$ matrix,
  $\tilde A = \left[\begin{smallmatrix} A & t \\ \mathbf{0} & 1 \end{smallmatrix}\right]$ with
  $A \in \mathbb{R}^{3\times3}$ and $t \in \mathbb{R}^3$. The $3 \times 3$ homogeneous matrix is
  the 2D case.

  (d) FALSE --- it is the other way around: $\det(J) > 1$ means local \emph{expansion},
  $\det(J) < 1$ local compression, and $\det(J) = 1$ no volume change. The divergence
  $\nabla \cdot T$, the trace of the Jacobian, carries related information.
\end{solution}

% =============================================================== Question 5 ===
\question
A free-form deformation (FFD) model with cubic B-splines is used for the non-linear registration
of two 3D images. The control points are placed on a uniform grid of $10 \times 12 \times 8$
points. How many parameters does this transformation model have, and how does this compare with
the number of parameters of a 3D affine transformation?

\begin{solution}
  Answer: $10 \cdot 12 \cdot 8 = 960$ control points, each carrying a displacement vector in
  $\mathbb{R}^3$, hence $960 \times 3 = 2880$ parameters. A 3D affine transformation has 12
  parameters (3 rotations $+$ 3 translations $+$ 3 scalings $+$ 3 shears).

  The FFD model therefore has roughly 240 times more degrees of freedom, which is exactly the
  point: the affine model is global and cannot capture local deformation, while the FFD model can.
  Note that the parameter count of the FFD depends only on the number of control points, not on
  the number of voxels --- that is what makes it tractable compared with one displacement vector
  per voxel.
\end{solution}
\answerspace{3.2cm}

% =============================================================== Question 6 ===
\question
The N3 algorithm corrects the bias field in MRI. It works in the log domain, where the observed
image is $j(x) = b(x) + i(x)$ with $b$ the bias field and $i$ the bias-free image, and where the
bias field is assumed to be independent of the image intensities. Each iteration alternates
between two estimation problems. Please briefly describe what each of them does, and explain why
the algorithm has to iterate with narrow Gaussians instead of correcting the bias in a single
step. You do not need to write down the equations.

\begin{solution}
  Answer: \emph{Estimating the distribution of the bias-free intensities.} Because $b$ and $i$
  are independent and add in the log domain, the observed intensity distribution is the
  convolution $J = I * B$. Since the multiplicative bias is close to 1, its logarithm is close to
  0, so $B$ can be approximated by a Gaussian with a small standard deviation and $J$ is a blurred
  version of $I$. The distribution $I$ is therefore recovered by \emph{sharpening} $J$, using a
  Wiener-style deblurring filter in the Fourier domain.

  \emph{Estimating the field.} Given $I$ and $B$, the expected bias-free intensity corresponding
  to an observed intensity $j$ is $E[i \mid j]$, so the bias estimate is $b_e(j) = j - E[i \mid j]$.
  This is a one-dimensional integral and is evaluated numerically. It is computed at anchor points
  and then smoothed over space with splines or radial basis functions --- this is where the prior
  knowledge that the bias field is spatially smooth enters --- and subtracted from $j$.

  Why iterate: $B$ is not actually known, it is only assumed to be Gaussian. Deconvolving with a
  wide Gaussian is unstable and can produce estimates that are not valid distributions at all.
  Since a wide Gaussian is a convolution of narrow ones, the correction is instead applied as a
  sequence of small, stable steps. Iterating also compensates for the other approximations made
  along the way: the noise term was ignored, and the deblurring filter is itself only approximate.
\end{solution}
\answerspace{7.5cm}

% =============================================================== Question 7 ===
\question
Which of the following statements on the loss metrics used for intensity-based registration is
\textbf{inaccurate}?

\begin{choices}
  \choice The sum of squared differences assumes that the intensities of the two images are
  directly comparable, which restricts it to intra-modality registration.
  \choice Normalized cross-correlation attains its maximum magnitude when the two images are
  related by a linear intensity relationship that is the same at every pixel; both positive and
  negative correlations indicate alignment, which is why the squared value is used in the loss.
  \choice Mutual information is maximized when the joint intensity distribution equals the product
  of the two marginals, that is, when the intensities of the two images are statistically
  independent.
  \choice Mutual information can capture more complicated, non-linear relationships between the
  intensities than normalized cross-correlation, but it is harder to compute.
\end{choices}

\begin{solution}
  Answer: C.

  It states exactly the condition under which mutual information is \emph{zero}:
  $\mathrm{MI} = 0$ if and only if $p(I,J) = p(I)p(J)$, which means the two images tell us nothing
  about each other. Registration \emph{maximizes} MI, because aligned images have a strong
  statistical relationship and their joint distribution is then as far as possible from the
  product of the marginals.
\end{solution}

% =============================================================== Question 8 ===
\question
Which of the following statements about interpretability and explainable AI (XAI) in medical image
analysis is \textbf{incorrect}?

\begin{choices}
  \choice Black-box, or model-agnostic, methods such as LIME require no access to the internals of
  the model, whereas gradient-based methods such as Grad-CAM are white-box methods.
  \choice Grad-CAM generalizes CAM: the importance weight of a feature map is the global average of
  the gradients of the class score with respect to that feature map, which makes it applicable to
  any CNN, while CAM requires a particular architecture.
  \choice Because saliency maps show which parts of the image the model actually uses, a model whose
  saliency lies largely outside the anatomical region of interest is necessarily performing poorly
  on the task.
  \choice TCAV explains a prediction in terms of human-defined concepts by measuring the directional
  derivative of the class score along a concept activation vector, and can be used to spot biases in
  a data set.
\end{choices}

\begin{solution}
  Answer: C.

  High accuracy and misplaced attention are entirely compatible --- that is precisely what makes
  shortcut learning dangerous. In the glioma grading case study, normalizing the intensities over
  the whole image instead of over the brain area \emph{raised} the AUC from 0.8857 to 0.9841, while
  the saliency maps showed the model attending outside the brain: the model was exploiting an
  artefact of the preprocessing. The pneumonia classifier that learned to recognize the clinical
  site from hospital-specific markers is the same story. Saliency outside the region of interest is
  a reason to distrust a good score, not evidence of a bad one.
\end{solution}

% =============================================================== Question 9 ===
\question
Which of the following statements about self-attention, as it was built up in the lecture, is
\textbf{incorrect}?

\begin{choices}
  \choice Parameter-free dot-product attention, $Y = \mathrm{Softmax}(XX^T)X$, has no capacity to
  learn from data, and every feature within a token plays an equal role in determining the attention
  coefficients.
  \choice Introducing a single learnable matrix $U$ and computing
  $Y = \mathrm{Softmax}(XUU^TX^T)\,XU$ still yields a symmetric matrix of attention scores, which is
  why separate key and query projections are introduced.
  \choice The scaling by $1/\sqrt{D_K}$ is used because the gradients of the softmax become very
  small for very large or very small inputs; if the entries of $Q$ and $K$ have unit variance, the
  entries of $QK^T$ have variance $D_K$.
  \choice In multi-head self-attention the outputs of the individual heads are averaged, so that the
  output dimension equals $D_V$ irrespective of the number of heads.
\end{choices}

\begin{solution}
  Answer: D.

  The heads are \emph{concatenated}, not averaged, and the concatenation is then projected back:
  $Y = \mathrm{Concat}(H_1, \dots, H_H)\,W_0$ with
  $\mathrm{Concat}(\cdot) \in \mathbb{R}^{N \times H D_V}$ and $W_0 \in \mathbb{R}^{H D_V \times D}$.
  Averaging would defeat the purpose of using several heads, which --- by analogy with the multiple
  filters of a convolutional layer --- is to let different heads attend to different patterns.
\end{solution}

% ============================================================== Question 10 ===
\question
In the Noise2Noise setting you have no clean images at all. For each input $x_n$ you only have
several \emph{noisy} versions $y_n^m$ of the underlying clean image $y_n$, and you train with
$\arg\min_\theta \sum_n \sum_m \|y_n^m - f(x_n;\theta)\|^2$. Under which condition does this give
the same solution as supervised training on clean targets, in the limit of infinite data?

\begin{choices}
  \choice When the noise is multiplicative rather than additive.
  \choice When the average of the noisy targets is the clean image, $\mathbb{E}[\hat y \mid x] = y$,
  because the prediction that minimizes an $\ell_2$ loss is the mean of the targets.
  \choice When at least half of the training targets are clean images.
  \choice When an adversarial loss is added to the $\ell_2$ loss, so that the outputs are forced to
  look realistic.
  \choice Never --- training on noisy targets always biases the solution towards the noise.
\end{choices}

\begin{solution}
  Answer: B.

  Minimizing $\mathbb{E}_{\hat y}[\|f(x;\theta) - \hat y\|^2]$ over the prediction gives
  $f(x;\theta) = \mathbb{E}_{\hat y}[\hat y]$: the $\ell_2$-optimal prediction is the mean of the
  targets it is trained against. So if the noisy targets are distributed around the clean image,
  $\mathbb{E}[\hat y \mid x] = y$, the network converges to the clean image even though it never
  sees one. The network cannot fit the individual noise realizations, because a single set of
  parameters has to satisfy all of them at once, and only their mean does.
\end{solution}

% ============================================================== Question 11 ===
\question
Consider a classification network that takes grey-scale input images of size $64 \times 64$ with
one input channel. All convolutions are \emph{valid} (no padding, $P = 0$), and the network uses
strided convolutions instead of pooling:

\begin{itemize}
  \item \textbf{Conv1}: 16 kernels with kernel size $K = 5 \times 5$ and stride $S = 2$ in each
    direction.
  \item \textbf{Conv2}: 32 kernels with $K = 5 \times 5$ and $S = 2$.
  \item \textbf{Conv3}: 64 kernels with $K = 3 \times 3$ and $S = 2$.
\end{itemize}

\noindent Recall the channel size of a strided, valid convolution:
$N_l = \lceil (N_{l-1} - k_1)/s_1 \rceil$. Circle True or False:

\begin{parts}
  \part \TFT\ \ The output of Conv3 has a spatial size of $5 \times 5$.

  \part \TFF\ \ The global receptive field of a neuron in Conv3 is $13 \times 13$.

  \part \TFF\ \ The global stride of a neuron in Conv3 is $6 \times 6$, since the three strides of
  2 add up.

  \part \TFT\ \ The total number of convolutional kernel weights (excluding the biases) is
  $5 \cdot 5 \cdot 1 \cdot 16 + 5 \cdot 5 \cdot 16 \cdot 32 + 3 \cdot 3 \cdot 32 \cdot 64
  = 31\,632$.
\end{parts}

\begin{solution}
  (a) TRUE --- applying $N_l = \lceil (N_{l-1} - k)/s \rceil$ layer by layer:
  Conv1 gives $\lceil (64-5)/2 \rceil = \lceil 29.5 \rceil = 30$;
  Conv2 gives $\lceil (30-5)/2 \rceil = \lceil 12.5 \rceil = 13$;
  Conv3 gives $\lceil (13-3)/2 \rceil = 5$.

  (b) FALSE --- it is $21 \times 21$. Tracking the receptive field $R$ and the jump $j$ from
  $R = 1$, $j = 1$ with $R \leftarrow R + (K-1)\cdot j$ and $j \leftarrow j \cdot S$:
  Conv1 $\rightarrow R = 5,\ j = 2$; Conv2 $\rightarrow R = 13,\ j = 4$;
  Conv3 $\rightarrow R = 21,\ j = 8$. The value $13$ is the receptive field one layer earlier,
  after Conv2.

  (c) FALSE --- strides compound multiplicatively: $2 \cdot 2 \cdot 2 = 8$, so the global stride
  is $8 \times 8$.

  (d) TRUE --- $400 + 12\,800 + 18\,432 = 31\,632$. Note that strided convolutions reduce the
  spatial dimensions quickly, but unlike max pooling they do \emph{not} buy any translation
  invariance, and they discard information --- the more so the larger the stride.
\end{solution}

% ============================================================== Question 12 ===
\question
Which of the following statements on restoration, synthesis and self-supervised learning are
correct (multiple answers are possible)?

\begin{choices}
  \choice Losses of the MSE type tend to produce over-smoothed restorations, which is one
  motivation for perceptual losses, where the two images are passed through a separately trained
  network and their intermediate features are compared.
  \choice In CycleGAN-style synthesis from unpaired data, the cycle-consistency term requires that
  mapping a source image to the target domain and back returns the original image.
  \choice In contrastive learning, the two views of an image have to be produced by geometric
  transformations only, since intensity changes would destroy the surrogate label.
  \choice A masked autoencoder uses a larger encoder on the visible parts of the image and a
  smaller decoder to reconstruct the masked parts; the trained encoder is then used as a backbone
  for downstream tasks.
  \choice Adversarial losses compare individual pairs of images, which is what makes them suitable
  for unpaired data sets.
\end{choices}

\begin{solution}
  Answer: (A), (B) and (D) are correct.

  (C) is wrong: the transformations used in contrastive learning are explicitly \emph{not}
  restricted to geometric ones and routinely include intensity changes --- the whole point is that
  the representation should be invariant to them. (The one place where transformations do have to
  be handled carefully is the \emph{local} contrastive loss, where they must not destroy the patch
  correspondences, so they are either restricted or undone before the patches are compared.)

  (E) is wrong in its premise, and the premise is what makes the conclusion work. An adversarial
  loss is a \emph{distributional} distance: the discriminator is not shown a pair of corresponding
  images, it is asked whether a single image looks like it came from the real set. That is exactly
  why it can be used when no correspondence between source and target images exists.
\end{solution}

\end{questions}

\end{document}
"
%% Compile the blank version (no answers):  pdflatex Sample_Exam_3.tex
%%
\ifx\withanswers\undefined
  \documentclass[11pt]{exam}
\else
  \documentclass[11pt,answers]{exam}
\fi

\usepackage[margin=2.4cm]{geometry}
\usepackage{amsmath,amssymb}
\usepackage{array}
\usepackage[T1]{fontenc}

% ---------------------------------------------------------------- page style
\pagestyle{foot}
\firstpageheader{}{}{}
\firstpagefooter{}{}{}
\runningheader{}{}{}
\runningfooter{}{}{Page \thepage}

% -------------------------------------------------------------- multiplechoice
% Reproduce the "(A) (B) (C) ..." labelling of the original exam.
\renewcommand{\choicelabel}{$\bigcirc$~(\Alph{choice})}

% ---------------------------------------------------------- true/false macros
% \TFT -> the statement is TRUE, \TFF -> the statement is FALSE.
% In the answer version the correct word is set in capitals and bold.
\newcommand{\TFT}{\ifprintanswers\textbf{TRUE}\hspace{1.1em}False\else True\hspace{1.1em}False\fi}
\newcommand{\TFF}{True\hspace{1.1em}\ifprintanswers\textbf{FALSE}\else False\fi}

% ------------------------------------------------------------- writing space
% Reserves room to write in the blank version; collapses in the answer version.
\newcommand{\answerspace}[1]{\ifprintanswers\else\vspace{#1}\fi}

\begin{document}

% ------------------------------------------------------------------ title page
\begin{center}
  {\LARGE\bfseries Sample Exam}\\[2mm]
  {\Large Medical Image Analysis}\\[6mm]
\end{center}

\vspace{4mm}
\noindent Student name:\hrulefill

\vspace{6mm}
\noindent Immatriculation number:\hrulefill

\vspace{10mm}

\begin{center}
  {\large\bfseries Exam Questions}
\end{center}

\noindent\textbf{Note:} All the questions have equal weight.

\vspace{4mm}

\begin{questions}

% =============================================================== Question 1 ===
\question
In 2D acquisitions, 2D image slices are acquired in one pre-defined direction one after the
other, and the final 3D volume is created by stacking the 2D slices. The voxel sizes are given
as in-plane (two values, the pixel size within the imaging plane) and through-plane (one value,
the slice thickness). Recall that the field-of-view is the real-world size of the entire image,
i.e.\ image size $\times$ voxel size. Based on this please compute the following.

\begin{parts}

  \part A \emph{coronal} acquisition with voxel size in-plane $0.9 \times 0.9$~mm$^2$ with 260
  pixels in both directions and voxel size through-plane $3.0$~mm with 64 pixels is acquired.
  What is the field-of-view in the head-to-toe, left-to-right and anterior-to-posterior
  directions, respectively?

  \begin{solution}
    Answer: $234 \times 234 \times 192$~mm$^3$.

    The coronal plane contains the head-to-toe and the left-to-right directions, so both in-plane
    directions give $260 \times 0.9 = 234$~mm. The through-plane direction of a coronal
    acquisition is anterior-to-posterior: $64 \times 3.0 = 192$~mm.
  \end{solution}
\answerspace{2.2cm}

  \part The protocol has to be changed so that the anterior-to-posterior coverage is at least
  210~mm, while the slice thickness stays at $3.0$~mm. How many slices have to be acquired?

  \begin{solution}
    Answer: 70 slices.

    $210 / 3.0 = 70$. The 64 slices of the original protocol only cover 192~mm.
  \end{solution}
\answerspace{1.8cm}

  \part You want to resample the volume of part (a) to isotropic voxels of
  $1.0 \times 1.0 \times 1.0$~mm$^3$. What are the resize factors and the resulting image size?
  In which of the three directions is the resampling least justified by the information actually
  contained in the data, and why?

  \begin{solution}
    Answer: resize factors $0.9 \times 0.9 \times 3.0$, resulting image size
    $234 \times 234 \times 192$~voxels.

    The resize factor per direction is the ratio of the old to the new voxel size: $0.9/1.0 = 0.9$
    head-to-toe, $0.9/1.0 = 0.9$ left-to-right and $3.0/1.0 = 3.0$ anterior-to-posterior.
    Applying these: $260 \cdot 0.9 = 234$, $260 \cdot 0.9 = 234$, $64 \cdot 3.0 = 192$. The
    field-of-view is unchanged, as it must be.

    The anterior-to-posterior direction. It is the through-plane direction, and it is being
    upsampled by a factor of 3: the interpolation invents two new samples between every pair of
    acquired slices. No new information is created --- the effective resolution in that direction
    is still that of a 3~mm slice, even though the voxel grid now says 1~mm.
  \end{solution}
\answerspace{4.5cm}

\end{parts}

% =============================================================== Question 2 ===
\question
Histogram equalization improves the contrast of an image by finding a monotonically increasing
intensity mapping $i \mapsto f(i)$ that flattens the histogram. Let $p(i)$ denote the probability
density function of the intensities, $P(i) = \int_0^i p(j)\,dj$ the cumulative distribution
function, and $i_{max}$ the maximum intensity. Which of the two mappings below is the one that
flattens the histogram? Please circle the correct choice.

\vspace{2mm}
\begin{center}
  $\displaystyle f(i) = i_{max}\, p(i)$
  \hspace{2.6cm}
  $\displaystyle f(i) = i_{max}\, P(i) = i_{max} \int_0^i p(j)\, dj$
\end{center}
\vspace{2mm}

\begin{solution}
  Answer: Right one --- the scaled \emph{cumulative} distribution function.

  Changing variables gives
  $p(f(i)) = p(i)\,\frac{di}{df(i)} = p(i) \cdot \frac{1}{p(i)} \cdot \frac{1}{i_{max}}
  = \frac{1}{i_{max}}$, i.e.\ a uniform (flat) density. Note also that only the CDF is guaranteed
  to be monotonically increasing, which the mapping has to be in order not to reorder intensities;
  the PDF is not.
\end{solution}
\answerspace{1.5cm}

% =============================================================== Question 3 ===
\question
When a Markov random field prior is used, segmentation can be written either as a posterior
maximization or as an energy minimization. $I$ is the image, $c$ the labelling, $G(x)$ the
neighbours of pixel $x$, and $Z$ a normalization constant.

\vspace{2mm}
\begin{center}
  $\displaystyle \max_c\ \log \Big[ \prod_{x \in \Omega} p(I(x) \mid c(x)) \Big]
   + \log \Big[ \tfrac{1}{Z} \exp\{-E(c)\} \Big]$
\end{center}
\vspace{1mm}
\begin{center}
  $\displaystyle \min_c\ \underbrace{\sum_{x \in \Omega} g\big(I(x) \mid \theta_{c(x)}\big)}_{\text{unary term}}
   \ +\ \underbrace{\sum_{x \in \Omega} \sum_{y \in G(x)} d\big(c(x), c(y)\big)}_{\text{pairwise term}}$
\end{center}
\vspace{2mm}

\begin{parts}
  \part Please indicate which term corresponds to which term in the two formulations.

  \begin{solution}
    Answer:
    \begin{align*}
      \log \textstyle\prod_{x} p(I(x) \mid c(x))
        &\ \Longleftrightarrow\ \text{the unary term}
        &&\text{(likelihood / data fidelity)}\\[1mm]
      \log \tfrac{1}{Z}\exp\{-E(c)\}
        &\ \Longleftrightarrow\ \text{the pairwise term}
        &&\text{(prior / neighbourhood consistency)}
    \end{align*}

    The correspondence requires the data model to be exponential, i.e.\
    $p(I(x) \mid c(x)) \propto \exp\{-g(I(x)\mid\theta_{c(x)})\}$, so that taking the logarithm of
    the product turns it into the unary sum. The pairwise term is exactly $E(c)$, and $\log Z$
    does not depend on $c$ and therefore drops out of the maximization.
  \end{solution}
\answerspace{2.8cm}

  \part Which theorem guarantees that a prior satisfying the Markov property can be written as a
  Gibbs distribution $\frac{1}{Z}\exp\{-E(c)\}$ with a locally defined energy $E$?

  \begin{solution}
    Answer: the Hammersley--Clifford theorem. It states that any probability distribution
    satisfying a Markovian property is a Gibbs distribution for an appropriate locally defined
    energy, and vice versa.
  \end{solution}
\answerspace{1.8cm}
\end{parts}

% =============================================================== Question 4 ===
\question
Please indicate whether the statement is true or false by circling the correct choice in each of
the following.

\begin{parts}
  \part \TFT\ \ A rigid-body transformation in 3D has 6 degrees of freedom: three rotation angles
  and three translations.

  \part \TFT\ \ For \emph{isotropic} scaling, $SRx = RSx$, i.e.\ the order in which rotation and
  scaling are applied does not matter.

  \part \TFF\ \ In homogeneous coordinates, a 3D affine transformation is written as a
  $3 \times 3$ matrix.

  \part \TFF\ \ The determinant of the Jacobian of a transformation is larger than 1 in the
  regions that the transformation compresses.
\end{parts}

\begin{solution}
  (a) TRUE --- three rotations, one about each axis, plus a translation vector in $\mathbb{R}^3$.
  In 2D the rigid model has 3 degrees of freedom (one angle, two translations).

  (b) TRUE --- with $S = sI$ the scaling commutes with any matrix. This is no longer the case for
  \emph{anisotropic} scaling, where $S = \mathrm{diag}(s_1, s_2, s_3)$: there the order matters a
  great deal, and applying the factors in the wrong order introduces shearing.

  (c) FALSE --- it is a $4 \times 4$ matrix,
  $\tilde A = \left[\begin{smallmatrix} A & t \\ \mathbf{0} & 1 \end{smallmatrix}\right]$ with
  $A \in \mathbb{R}^{3\times3}$ and $t \in \mathbb{R}^3$. The $3 \times 3$ homogeneous matrix is
  the 2D case.

  (d) FALSE --- it is the other way around: $\det(J) > 1$ means local \emph{expansion},
  $\det(J) < 1$ local compression, and $\det(J) = 1$ no volume change. The divergence
  $\nabla \cdot T$, the trace of the Jacobian, carries related information.
\end{solution}

% =============================================================== Question 5 ===
\question
A free-form deformation (FFD) model with cubic B-splines is used for the non-linear registration
of two 3D images. The control points are placed on a uniform grid of $10 \times 12 \times 8$
points. How many parameters does this transformation model have, and how does this compare with
the number of parameters of a 3D affine transformation?

\begin{solution}
  Answer: $10 \cdot 12 \cdot 8 = 960$ control points, each carrying a displacement vector in
  $\mathbb{R}^3$, hence $960 \times 3 = 2880$ parameters. A 3D affine transformation has 12
  parameters (3 rotations $+$ 3 translations $+$ 3 scalings $+$ 3 shears).

  The FFD model therefore has roughly 240 times more degrees of freedom, which is exactly the
  point: the affine model is global and cannot capture local deformation, while the FFD model can.
  Note that the parameter count of the FFD depends only on the number of control points, not on
  the number of voxels --- that is what makes it tractable compared with one displacement vector
  per voxel.
\end{solution}
\answerspace{3.2cm}

% =============================================================== Question 6 ===
\question
The N3 algorithm corrects the bias field in MRI. It works in the log domain, where the observed
image is $j(x) = b(x) + i(x)$ with $b$ the bias field and $i$ the bias-free image, and where the
bias field is assumed to be independent of the image intensities. Each iteration alternates
between two estimation problems. Please briefly describe what each of them does, and explain why
the algorithm has to iterate with narrow Gaussians instead of correcting the bias in a single
step. You do not need to write down the equations.

\begin{solution}
  Answer: \emph{Estimating the distribution of the bias-free intensities.} Because $b$ and $i$
  are independent and add in the log domain, the observed intensity distribution is the
  convolution $J = I * B$. Since the multiplicative bias is close to 1, its logarithm is close to
  0, so $B$ can be approximated by a Gaussian with a small standard deviation and $J$ is a blurred
  version of $I$. The distribution $I$ is therefore recovered by \emph{sharpening} $J$, using a
  Wiener-style deblurring filter in the Fourier domain.

  \emph{Estimating the field.} Given $I$ and $B$, the expected bias-free intensity corresponding
  to an observed intensity $j$ is $E[i \mid j]$, so the bias estimate is $b_e(j) = j - E[i \mid j]$.
  This is a one-dimensional integral and is evaluated numerically. It is computed at anchor points
  and then smoothed over space with splines or radial basis functions --- this is where the prior
  knowledge that the bias field is spatially smooth enters --- and subtracted from $j$.

  Why iterate: $B$ is not actually known, it is only assumed to be Gaussian. Deconvolving with a
  wide Gaussian is unstable and can produce estimates that are not valid distributions at all.
  Since a wide Gaussian is a convolution of narrow ones, the correction is instead applied as a
  sequence of small, stable steps. Iterating also compensates for the other approximations made
  along the way: the noise term was ignored, and the deblurring filter is itself only approximate.
\end{solution}
\answerspace{7.5cm}

% =============================================================== Question 7 ===
\question
Which of the following statements on the loss metrics used for intensity-based registration is
\textbf{inaccurate}?

\begin{choices}
  \choice The sum of squared differences assumes that the intensities of the two images are
  directly comparable, which restricts it to intra-modality registration.
  \choice Normalized cross-correlation attains its maximum magnitude when the two images are
  related by a linear intensity relationship that is the same at every pixel; both positive and
  negative correlations indicate alignment, which is why the squared value is used in the loss.
  \choice Mutual information is maximized when the joint intensity distribution equals the product
  of the two marginals, that is, when the intensities of the two images are statistically
  independent.
  \choice Mutual information can capture more complicated, non-linear relationships between the
  intensities than normalized cross-correlation, but it is harder to compute.
\end{choices}

\begin{solution}
  Answer: C.

  It states exactly the condition under which mutual information is \emph{zero}:
  $\mathrm{MI} = 0$ if and only if $p(I,J) = p(I)p(J)$, which means the two images tell us nothing
  about each other. Registration \emph{maximizes} MI, because aligned images have a strong
  statistical relationship and their joint distribution is then as far as possible from the
  product of the marginals.
\end{solution}

% =============================================================== Question 8 ===
\question
Which of the following statements about interpretability and explainable AI (XAI) in medical image
analysis is \textbf{incorrect}?

\begin{choices}
  \choice Black-box, or model-agnostic, methods such as LIME require no access to the internals of
  the model, whereas gradient-based methods such as Grad-CAM are white-box methods.
  \choice Grad-CAM generalizes CAM: the importance weight of a feature map is the global average of
  the gradients of the class score with respect to that feature map, which makes it applicable to
  any CNN, while CAM requires a particular architecture.
  \choice Because saliency maps show which parts of the image the model actually uses, a model whose
  saliency lies largely outside the anatomical region of interest is necessarily performing poorly
  on the task.
  \choice TCAV explains a prediction in terms of human-defined concepts by measuring the directional
  derivative of the class score along a concept activation vector, and can be used to spot biases in
  a data set.
\end{choices}

\begin{solution}
  Answer: C.

  High accuracy and misplaced attention are entirely compatible --- that is precisely what makes
  shortcut learning dangerous. In the glioma grading case study, normalizing the intensities over
  the whole image instead of over the brain area \emph{raised} the AUC from 0.8857 to 0.9841, while
  the saliency maps showed the model attending outside the brain: the model was exploiting an
  artefact of the preprocessing. The pneumonia classifier that learned to recognize the clinical
  site from hospital-specific markers is the same story. Saliency outside the region of interest is
  a reason to distrust a good score, not evidence of a bad one.
\end{solution}

% =============================================================== Question 9 ===
\question
Which of the following statements about self-attention, as it was built up in the lecture, is
\textbf{incorrect}?

\begin{choices}
  \choice Parameter-free dot-product attention, $Y = \mathrm{Softmax}(XX^T)X$, has no capacity to
  learn from data, and every feature within a token plays an equal role in determining the attention
  coefficients.
  \choice Introducing a single learnable matrix $U$ and computing
  $Y = \mathrm{Softmax}(XUU^TX^T)\,XU$ still yields a symmetric matrix of attention scores, which is
  why separate key and query projections are introduced.
  \choice The scaling by $1/\sqrt{D_K}$ is used because the gradients of the softmax become very
  small for very large or very small inputs; if the entries of $Q$ and $K$ have unit variance, the
  entries of $QK^T$ have variance $D_K$.
  \choice In multi-head self-attention the outputs of the individual heads are averaged, so that the
  output dimension equals $D_V$ irrespective of the number of heads.
\end{choices}

\begin{solution}
  Answer: D.

  The heads are \emph{concatenated}, not averaged, and the concatenation is then projected back:
  $Y = \mathrm{Concat}(H_1, \dots, H_H)\,W_0$ with
  $\mathrm{Concat}(\cdot) \in \mathbb{R}^{N \times H D_V}$ and $W_0 \in \mathbb{R}^{H D_V \times D}$.
  Averaging would defeat the purpose of using several heads, which --- by analogy with the multiple
  filters of a convolutional layer --- is to let different heads attend to different patterns.
\end{solution}

% ============================================================== Question 10 ===
\question
In the Noise2Noise setting you have no clean images at all. For each input $x_n$ you only have
several \emph{noisy} versions $y_n^m$ of the underlying clean image $y_n$, and you train with
$\arg\min_\theta \sum_n \sum_m \|y_n^m - f(x_n;\theta)\|^2$. Under which condition does this give
the same solution as supervised training on clean targets, in the limit of infinite data?

\begin{choices}
  \choice When the noise is multiplicative rather than additive.
  \choice When the average of the noisy targets is the clean image, $\mathbb{E}[\hat y \mid x] = y$,
  because the prediction that minimizes an $\ell_2$ loss is the mean of the targets.
  \choice When at least half of the training targets are clean images.
  \choice When an adversarial loss is added to the $\ell_2$ loss, so that the outputs are forced to
  look realistic.
  \choice Never --- training on noisy targets always biases the solution towards the noise.
\end{choices}

\begin{solution}
  Answer: B.

  Minimizing $\mathbb{E}_{\hat y}[\|f(x;\theta) - \hat y\|^2]$ over the prediction gives
  $f(x;\theta) = \mathbb{E}_{\hat y}[\hat y]$: the $\ell_2$-optimal prediction is the mean of the
  targets it is trained against. So if the noisy targets are distributed around the clean image,
  $\mathbb{E}[\hat y \mid x] = y$, the network converges to the clean image even though it never
  sees one. The network cannot fit the individual noise realizations, because a single set of
  parameters has to satisfy all of them at once, and only their mean does.
\end{solution}

% ============================================================== Question 11 ===
\question
Consider a classification network that takes grey-scale input images of size $64 \times 64$ with
one input channel. All convolutions are \emph{valid} (no padding, $P = 0$), and the network uses
strided convolutions instead of pooling:

\begin{itemize}
  \item \textbf{Conv1}: 16 kernels with kernel size $K = 5 \times 5$ and stride $S = 2$ in each
    direction.
  \item \textbf{Conv2}: 32 kernels with $K = 5 \times 5$ and $S = 2$.
  \item \textbf{Conv3}: 64 kernels with $K = 3 \times 3$ and $S = 2$.
\end{itemize}

\noindent Recall the channel size of a strided, valid convolution:
$N_l = \lceil (N_{l-1} - k_1)/s_1 \rceil$. Circle True or False:

\begin{parts}
  \part \TFT\ \ The output of Conv3 has a spatial size of $5 \times 5$.

  \part \TFF\ \ The global receptive field of a neuron in Conv3 is $13 \times 13$.

  \part \TFF\ \ The global stride of a neuron in Conv3 is $6 \times 6$, since the three strides of
  2 add up.

  \part \TFT\ \ The total number of convolutional kernel weights (excluding the biases) is
  $5 \cdot 5 \cdot 1 \cdot 16 + 5 \cdot 5 \cdot 16 \cdot 32 + 3 \cdot 3 \cdot 32 \cdot 64
  = 31\,632$.
\end{parts}

\begin{solution}
  (a) TRUE --- applying $N_l = \lceil (N_{l-1} - k)/s \rceil$ layer by layer:
  Conv1 gives $\lceil (64-5)/2 \rceil = \lceil 29.5 \rceil = 30$;
  Conv2 gives $\lceil (30-5)/2 \rceil = \lceil 12.5 \rceil = 13$;
  Conv3 gives $\lceil (13-3)/2 \rceil = 5$.

  (b) FALSE --- it is $21 \times 21$. Tracking the receptive field $R$ and the jump $j$ from
  $R = 1$, $j = 1$ with $R \leftarrow R + (K-1)\cdot j$ and $j \leftarrow j \cdot S$:
  Conv1 $\rightarrow R = 5,\ j = 2$; Conv2 $\rightarrow R = 13,\ j = 4$;
  Conv3 $\rightarrow R = 21,\ j = 8$. The value $13$ is the receptive field one layer earlier,
  after Conv2.

  (c) FALSE --- strides compound multiplicatively: $2 \cdot 2 \cdot 2 = 8$, so the global stride
  is $8 \times 8$.

  (d) TRUE --- $400 + 12\,800 + 18\,432 = 31\,632$. Note that strided convolutions reduce the
  spatial dimensions quickly, but unlike max pooling they do \emph{not} buy any translation
  invariance, and they discard information --- the more so the larger the stride.
\end{solution}

% ============================================================== Question 12 ===
\question
Which of the following statements on restoration, synthesis and self-supervised learning are
correct (multiple answers are possible)?

\begin{choices}
  \choice Losses of the MSE type tend to produce over-smoothed restorations, which is one
  motivation for perceptual losses, where the two images are passed through a separately trained
  network and their intermediate features are compared.
  \choice In CycleGAN-style synthesis from unpaired data, the cycle-consistency term requires that
  mapping a source image to the target domain and back returns the original image.
  \choice In contrastive learning, the two views of an image have to be produced by geometric
  transformations only, since intensity changes would destroy the surrogate label.
  \choice A masked autoencoder uses a larger encoder on the visible parts of the image and a
  smaller decoder to reconstruct the masked parts; the trained encoder is then used as a backbone
  for downstream tasks.
  \choice Adversarial losses compare individual pairs of images, which is what makes them suitable
  for unpaired data sets.
\end{choices}

\begin{solution}
  Answer: (A), (B) and (D) are correct.

  (C) is wrong: the transformations used in contrastive learning are explicitly \emph{not}
  restricted to geometric ones and routinely include intensity changes --- the whole point is that
  the representation should be invariant to them. (The one place where transformations do have to
  be handled carefully is the \emph{local} contrastive loss, where they must not destroy the patch
  correspondences, so they are either restricted or undone before the patches are compared.)

  (E) is wrong in its premise, and the premise is what makes the conclusion work. An adversarial
  loss is a \emph{distributional} distance: the discriminator is not shown a pair of corresponding
  images, it is asked whether a single image looks like it came from the real set. That is exactly
  why it can be used when no correspondence between source and target images exists.
\end{solution}

\end{questions}

\end{document}
